\chapter{Truncations of the beta lattice and Hermite coefficients}\label{chap:tr}\chaptermark{Truncations and Hermite coefficients}

Section~\ref{ch:rz4} truncates the lattice of Biane, Pitman and Yor for zeta after $n$ terms. The truncated Mellin transforms are explicit, their zeros move to the zeros of $\xi$ at the first order rate $\kappa(\rho)/n$ (Proposition~\ref{rz:prop:firstorder}), the exact form of $\kappa(\rho)$ ties the side of approach to the Riemann Siegel phase and so to Gram's law (Theorem~\ref{rz:thm:kappaexact}), and a single zero approached from the wrong side, $\rho_{213}$, shows that the one sided conjecture is false (Corollary~\ref{rz:cor:Tfalse}). Symmetrising by the functional equation removes the mechanism and leaves a statement that is not refuted (Proposition~\ref{rz:prop:symtrunc}). This chapter carries that programme to the Dirichlet beta function $L(s,\chi_4)$, where the lattice is $\lambda_k=2\pi^2(k-\frac12)^2$ and the Mellin transform is the completed $L$ function, and it adds what the zeta section does not have: the orders $m\ge1$ of cumulant matching, an exact mean corrected truncation as a finite sum of incomplete gamma functions, and a computation of the drift at all $379$ zeros below height $500$. The second half is a separate subject, the expansion of $\Xi$ in Hermite functions at a small scale, whose coefficients are samples of $\Xi$ at the turning points of the Hermite functions, up to an amplitude factor and an error $O(a^4)$, and have been used to count zeros from the P\'olya kernel alone.

The chapter develops two notes of Nick Polson, on the Thorin truncation of $L(s,\chi_4)$ \cite{P26trunc} and on progress with open problems about the Hermite coefficients of $\Xi$ \cite{P26hermite}, together with the computations done for this book (the scripts are \texttt{code/tr\_*.py}, listed in \texttt{code/README.md}). Where the recomputation differed from a note the recomputed value is printed, and a statement that was not recomputed is marked unchecked. Several statements of the notes were restated in the course of the merge: the drift law of the first note holds for a remainder only under a hypothesis on its first unmatched cumulant (Theorem~\ref{tr:thm:drift}, Proposition~\ref{tr:prop:gamma}), the corollary that no one sided criterion can hold is conditional on one zero with a negative drift (Corollary~\ref{tr:cor:nogo}), and the asymptotic of the second note is proved in an additive form that is uniform in $m$ and fails in the relative form at the zeros of $\Xi$ (Theorem~\ref{tr:thm:herm}, Proposition~\ref{tr:prop:zeros}).

Notation is the book's, with these warnings. The function $\xi_4(s)=\xi(s,\chi_4)$ is that of Section~\ref{sec:dbeta:beta}, and the letter $\Lambda$, which the first note uses for $\frac{\sqrt\pi}2\xi_4$, is the de Bruijn Newman constant here. The variable $\omega$ is the exit time of Remark~\ref{rem:dbeta:exit} divided by $2\pi$. The functions $\vartheta_\beta$ and $Z_\beta$ are the Riemann Siegel and Hardy functions of $\chi_4$, written with a subscript so that $\vartheta$ and $Z$ remain those of zeta. The kernel $\Phi$ is that of \eqref{eq:ch2:Phi} (the second note does not define it), $H_t$ is the de Bruijn Newman function of Section~\ref{ch:dbn}, and $H_n$ with an integer index is the Hermite polynomial of Section~\ref{ht:sec:jensen}. The letter $a$ is the Hermite scale in the second half and a constant is written $\lambda_*=2\pi^2$ in the first.

Status words follow the rest of the book. A statement is proved when a proof or a pointer is given, computed when a script of the book produced it, conditional, heuristic or open as stated. A computation in floating point is called computed and never certified; the zero counts of this chapter are of that kind. The results are these.
\begin{enumerate}
\item GRH for $\chi_4$ is equivalent to a local zero free condition on the truncated Mellin transforms (Theorem~\ref{tr:thm:local}, proved). The truncated zeros move by $\delta_mn^{-(2m+1)}$ for the hard cutoff and the mean correction, and by a multiple of the first unmatched cumulant difference in general (Theorem~\ref{tr:thm:drift}, Corollary~\ref{tr:cor:named}, proved); the formula for $m\ge2$ with a generalised gamma convolution remainder is false (Proposition~\ref{tr:prop:gamma}, proved), and the drift of the first zero is computed for the three kinds of remainder (Table~\ref{tr:tab:drift}).
\item The sign of the real part of the drift is a Gram parity of $\chi_4$ with two corrections (Theorem~\ref{tr:thm:phase}, proved). Among the $379$ zeros below height $500$ the mean correction approaches from the left exactly at the five zeros with anomalous Gram parity (computed, Table~\ref{tr:tab:gram}). A zero with negative drift excludes a one sided criterion at that order (Corollary~\ref{tr:cor:nogo}, proved), and the premise is computed and not proved.
\item The truncations have transient zeros in the left half strip (Table~\ref{tr:tab:counts}, computed for the hard cutoff and the mean correction). The symmetrised truncations satisfy a Hurwitz criterion (Theorem~\ref{tr:thm:symhurwitz}, proved) and keep every simple critical zero on the line locally (Proposition~\ref{tr:prop:symlocal}, proved); the heights reached are those of the computed tables, and the mean corrected rows are computed as well.
\item The Hermite coefficients of $\Xi$ at scale $a$ are the Taylor coefficients of $e^{z^2/2}$ times the backward heat flow of $\Xi$ for the time $2a^2$ (Theorem~\ref{tr:thm:heat}, proved). They sample $\Xi$ at the turning points $a\sqrt{4m+1}$ with an additive error $O(a^4)$ that is uniform in $m$ (Theorem~\ref{tr:thm:herm}, proved), and the first negative coefficient reads off $\gamma_1$ (Corollary~\ref{tr:cor:first}, proved; computed, Table~\ref{tr:tab:first}).
\item The number of sign changes of the coefficients equals the number of zeros of zeta up to height $1000$ at scales up to about $5$ to $6$, with a lag of order one at the top edge (Remark~\ref{tr:obs:count}, computed); what a verified count would prove is Proposition~\ref{tr:prop:verified}.
\end{enumerate}
The open problems close each half.
